\chapter{Backgrounds}

\section{Software Contracts}

\subsection{Overview}

\subsection{Formal Language}

In this section we define the formal language with support for software contracts. This language is influenced by \cite{cpcf}, with indy semantics for dependent contracts defined in \cite{completemonitors}. We name this language Contracts.

The surface syntax for Contracts language is provided in Table \ref{tab:contracts_syntax}. The type rules are provided in Table \ref{tab:contracts_type_rules}. Reductions are provided in Table \ref{tab:contracts_reduction}.

\begin{singlespace}
\begin{itemize}
    \item[] \begin{tabular}{llll}
        Types & $\tau$ & $::=$ & $\Tnum \mid \Tbool \mid \Tarrow{\tau_1}{\tau_2} \mid \Tcon{\tau}$ \\
        Contracts & $\kappa$ & $::=$ & $\Kflat{e} \mid \Karrow{\kappa_1}{\kappa_2} \mid \kappa_1 \xrightarrow{d} (\lambda x:\tau. \kappa_2)$ \\
        Expressions & $e$ & $::=$ & $0, 1, ... \mid \ttval \mid \ffval \mid x \mid \lam{x}{\tau}{e} \mid \fix{x}{\tau}{e}$ \\
        & & $\mid$ & $e_1 + e_2 \mid e_1 \wedge e_2 \mid e_1 \vee e_2 \mid e_1 \ e_2$ \\
        & & $\mid$ & $\text{if } e_1 \text{ then } e_2 \text{ else } e_3 \mid \zeroq{e}$ \\
        & & $\mid$ & $\mon{k}{l}{j}{\kappa}{e} \mid \error{l}{p}$ \\
        Values & $v$ & $::=$ & $0, 1, ... \mid \ttval \mid \ffval \mid λ x. e \mid \mu x. e$ \\
        Contexts & $E$ & $::=$ & $\hole \mid E + e \mid v + E \mid E \wedge e \mid v \wedge E \mid E \vee e \mid v \vee E$ \\
        & & $\mid$ & $E \ e \mid v \ E \mid \mon{k}{l}{j}{\kappa}{E}$
    \end{tabular}
\end{itemize}

\begin{mathpar}
\inferrule[If-T]
  {\strut}
  {\text{if } \ttval \text{ then } e_1 \text{ else } e_2 \red e_1}

\inferrule[If-F]
  {\strut}
  {\text{if } \ffval \text{ then } e_1 \text{ else } e_2 \red e_2}

\inferrule[Mon-Flat]
  {\strut}
  {\mon{k}{l}{j}{\Kflat{e}}{v} \red \text{if } e \ v \text{ then } v \text{ else } \error{j}{k}}

\inferrule[Mon-Arr]
  {\strut}
  {\mon{k}{l}{j}{\Karrow{\kappa_1}{\kappa_2}}{v} \red
   λ x . \mon{k}{l}{j}{\kappa_2}{v \ (\mon{l}{k}{j}{\kappa_1}{x})}}

\inferrule[Mon-Dep]
  {\strut}
  {\mon{k}{l}{j}{\kappa_1 \xrightarrow{d} (\lam{x}{\tau}{\kappa_2})}{v} \red
   λ x. \mon{k}{l}{j}{\Ksubst{\mon{l}{j}{j}{\kappa_1}{x}}{x}{\kappa_2}}{v \ (\mon{l}{k}{j}{\kappa_1}{x})}}
\end{mathpar}
\end{singlespace}

\begin{singlespace}
\begin{mathpar}
\inferrule[Context]
  {e_1 \red e_2}
  {E[e_1] \step E[e_2]}

\inferrule[Err-Prop]
  {\strut}
  {E[\error{l}{p}] \step \error{l}{p}}
\end{mathpar}

\begin{mathpar}
\inferrule[T-Error]
  {\strut}
  {\types{\ctx}{\error{l}{p}}{\tau}}

\inferrule[T-Mon]
  {\types{\ctx}{\kappa}{\Tcon{\tau}} \\ \types{\ctx}{e}{\tau}}
  {\types{\ctx}{\mon{k}{l}{j}{\kappa}{e}}{\tau}}

\inferrule[T-Flat]
  {\types{\ctx}{e}{\Tarrow{o}{\mathsf{bool}}}}
  {\types{\ctx}{\Kflat{e}}{\Tcon{o}}}

\inferrule[T-Arr]
  {\types{\ctx}{\kappa_1}{\Tcon{\tau_1}} \\ \types{\ctx}{\kappa_2}{\Tcon{\tau_2}}}
  {\types{\ctx}{\Karrow{\kappa_1}{\kappa_2}}{\Tcon{\Tarrow{\tau_1}{\tau_2}}}}

\inferrule[T-Dep-Arr]
  {\types{\ctx}{\kappa_1}{\Tcon{\tau_1}} \\ \types{\ext{\ctx}{x}{\tau_1}}{\kappa_2}{\Tcon{\tau_2}}}
  {\types{\ctx}{\Kdarrow{\kappa_1}{x}{\tau_1}{\kappa_2}}{\Tcon{\Tarrow{\tau_1}{\tau_2}}}}
\end{mathpar}

\end{singlespace}


Following the paper Effectful Contracts for the proofs needed.

We define a partial evaluation function $\kw{eval}$ that takes a closed expression and yields an answer. ...

\begin{theorem}[Functional Evaluation]
  Two facts about the evaluator holds:

  \begin{itemize}
    \item[] $\kw{eval}$ is a partial function, or
    \item[] If $e$ is a program, then either (i) $\kw{eval}(e)$ is defined or (ii) reduction starting from $e$ is unbounded
  \end{itemize}

\end{theorem}

Define good value $b=?$.

Contracts enhances a program by placing monitors to watch for violations. With this premise, a program after placing monitors should maintain its semantics, or signal an error. Such property is \textit{contract erasure}. We define an erasure function $\mathscr{E}$ removing all monitors placed on the program.

\begin{theorem}[Contract Erasure]
  If $\kw{eval}(e) = b$ then $\kw{eval}(\mathscr{E}(e)) = b$

\end{theorem}


Finally, the Contracts language should satisfies \textit{blame correctness}. In other words, if the program signals an error, it should blame the correct component.

\begin{theorem}[Blame Correctness]

\end{theorem}


\section{Effect Handlers}

\subsection{Overview}

\subsection{Formal Language}

In this section, we define the formal language with support for effect handlers. This language is influenced by \cite{introeffecthandlers}. We name this language Effects.

The type system in Effects separates between a basic value and a computation. A computation type is accompanied with a list of possible effects, denoting a series of effects that are possible when running the computation. Many expression is of computation type, therefore to use a value, one may need to wrap it into a computation using $\kw{return}$.
An effect $\mathsf{op}$ is invoked with argument $v$ through $\vop{op}{v}{y.c}$. Once the result is available, it will be available as $y$ in the continuation $c$. The continuation is helpful to keep reduction and bubbling up effect invocation to a handler that has it.

% In constrast, \cite{leijen2016algebraic} uses two evaluation contexts

\begin{singlespace}

\begin{itemize}
    \item[] \begin{tabular}{llll}
        Value $v$ & $::=$ & $x \mid \vtrue \mid \vfalse \mid n \mid \vfun{x}{c} \mid h$ & \text{Values} \\
        Handler $h$ & $::=$ & $\vhandler \{ \kw{return} \ x \mapsto c_r, \mathsf{op}_i(x; k) \mapsto c_i \}$ & \text{Handlers} \\
        Computation $c$ & $::=$ & $\vreturn{v} \mid \vop{op}{v}{y.c} \mid \vdo{x}{c_1}{c_2}$ & \text{Computations} \\
        & $\mid$ & $\vif{v}{c_1}{c_2} \mid v_1 \ v_2 \mid \vwith{v}{c}$ & \\
        Value Type $A, B$ & $::=$ & $\Tnum \mid \Tbool \mid A \to \underline{C} \mid \underline{C} \Rightarrow \underline{D}$ & \\
        Comp. Type $\underline{C}, \underline{D}$ & $::=$ & $A!\{\mathsf{op}_1, \dots, \mathsf{op}_n\}$ & \\
    \end{tabular}
\end{itemize}

\begin{mathpar}
\inferrule[Do-Ret]
  {\strut}
  {\vdo{x}{\vreturn{v}}{c} \effstep \Ksubst{v}{x}{c}}

\inferrule[Do-Op]
  {\strut}
  {\vdo{x}{\vop{op}{v}{y.c_1}}{c_2} \effstep \vop{op}{v}{y. \vdo{x}{c_1}{c_2}}}

\inferrule[App]
  {\strut}
  {(\vfun{x}{c}) \ v \effstep \Ksubst{v}{x}{c}}

\inferrule[With-Ret]
  {h = \vhandler \{ \kw{return} \ x \mapsto c_r, \dots \}}
  {\vwith{h}{\vreturn{v}} \effstep \Ksubst{v}{x}{c_r}}

\inferrule[With-Op]
  {h = \vhandler \{ \dots, \mathsf{op}_i(x; k) \mapsto c_i, \dots \}}
  {\vwith{h}{\vop{op_i}{v}{y.c} \effstep \Ksubst{v}{x}{\Ksubst{\vfun{y}{\vwith{h}{c}}}{k}{c_i}}}}

\inferrule[With-Prop]
  {\mathsf{op} \notin \{\mathsf{op}_1, \dots, \mathsf{op}_n\}}
  {\vwith{h}{\vop{op}{v}{y.c}} \effstep \vop{op}{v}{y. \vwith{h}{c}}}
\end{mathpar}

\begin{mathpar}
\inferrule[T-Handler]
  {\types{\ext{\ctx}{x}{A}}{c_r}{B!\Delta'} \\
   [\types{\Gamma, x:A_i, k:\Tarrow{B_i}{B!\Delta'}}{c_i}{B!\Delta'}]_{1 \leq i \leq n} \\
   \Delta \setminus \{\mathsf{op}_1, \dots, \mathsf{op}_n\} \subseteq \Delta'}
  {\types{\ctx}{\vhandler \{ \vreturn x \mapsto c_r, \mathsf{op}_i(x; k) \mapsto c_i, \dots \mathsf{op}_n(x; k) \mapsto c_n\}}
    {A!Δ  \Rightarrow B!Δ'}}

\inferrule[T-Op]
  {(\mathsf{op} : A_{\mathsf{op}} \to B_{\mathsf{op}}) \in \Sigma \\ \types{\ctx}{v}{A_{\mathsf{op}}} \\
   \types{\ext{\ctx}{y}{B_{\mathsf{op}}}}{c}{A!\Delta} \\ \mathsf{op} \in \Delta}
  {\types{\ctx}{\vop{op}{v}{y.c}}{A!\Delta}}

\inferrule[T-With]
  {\types{\ctx}{v}{\underline{C} \Rightarrow \underline{D}} \\ \types{\ctx}{c}{\underline{C}}}
  {\types{\ctx}{\vwith{v}{c}}{\underline{D}}}
\end{mathpar}

\end{singlespace}

\begin{theorem}
Effects is safe by Progress and Preservation. If $\types{}{c}{A!Δ}$ then:

\begin{itemize}
  \item[] $c = \vreturn v$ for some $\types{}{v}{A}$, or
  \item[] $c = \vop{op}{v}{y.c}$ for some $\mathsf{op} ∈ Δ$, or
  \item[] Exists $\types{}{c'}{A!Δ}$ such that $c \rightarrow c'$
\end{itemize}

\end{theorem}




\section{OCaml}

Yapping about OCaml here.

\subsection{Effect Handlers support}

OCaml added support for Effect Handlers since OCaml 5.0, released in 2022. The feature was added to Multicore OCaml following the work of \cite{sivaramakrishnan2021retrofitting}, which is later merged into main branch of development. Before the official support of Effect Handlers, \cite{kiselyov2018eff} implemented Effect Handlers by embedding Eff into OCaml using Module Functors.

\subsection{Software Contracts support}

Software Contracts are not realized in OCaml. Previously, \cite{xu2014dynamic} implemented software contracts into OCaml. However, it has never been merged into the official OCaml.

\subsection{Metaprogramming in OCaml}

OCaml has support for metaprogramming through preprocessing. A preprocessor can modify the AST freely before compilation stage.
